\chapter{Theoretical Foundations and Literature Review}
\label{ch:theory}

This chapter establishes the conceptual foundations on which the thesis rests and positions its
contribution within existing research. It proceeds from the operational context
(\Cref{sec:lit-warehouse}) through the two decisions under study --- capacity
(\Cref{sec:lit-capacity}) and sequencing (\Cref{sec:lit-sequencing}) --- to the methodological
apparatus used to study them (\Cref{sec:lit-des,sec:lit-experimentation}), then to learning-based
control (\Cref{sec:lit-rl}) and to the requirement that a planning instrument be usable
(\Cref{sec:lit-dss}). \Cref{sec:lit-gap} states the position this thesis occupies.

The intention is to build an argument rather than to catalogue studies. Where a body of work is
summarised, it is summarised because it supports a step in that argument.

\section{Warehouse fulfilment operations}
\label{sec:lit-warehouse}

\subsection{The operational context}

A warehouse converts a stream of customer orders into a stream of dispatched shipments. Research
on how it should do so has been organised for some decades around a hierarchy of decisions:
strategic questions of design and layout, tactical questions of resource and storage allocation,
and operational questions of how work is executed day to day \citep{gu2007,gu2010}. The
comprehensive reviews in that tradition make clear that the literature is unevenly distributed
across the hierarchy, with a heavy concentration on design and on the execution of picking.

That concentration is well justified economically. Order picking is repeatedly identified as the
most labour-intensive activity in a typical warehouse and the largest single component of
operating expenditure, with estimates placing it at over half of total warehouse operating cost
\citep{dekoster2007}. If one activity dominates cost, improving it dominates the return to
research effort.

The resulting body of work is largely concerned with how the picking process itself should be
organised. Comparative studies of storage assignment and picker routing establish substantial
achievable savings in travel and handling effort \citep{petersen2004}, and a parallel literature
addresses the design and control of automated storage and retrieval systems as an alternative to
manual picking altogether \citep{roodbergen2009}. Both strands take the volume of work as given
and ask how to execute it more efficiently. The present study takes the execution technology as
given and asks how much labour to buy and in what order to release the work --- a different
question at a different planning horizon, and one that remains relevant regardless of how
efficient the underlying picking process is.

\subsection{What e-commerce changed}

The character of the problem shifted materially with the growth of online retail.
\citet{boysen2019} survey the consequences and identify a consistent pattern: order sizes fall,
order frequency rises, demand becomes more volatile, and --- most consequentially for this
thesis --- delivery promises shorten and diversify. A warehouse serving a retail supply chain
handled comparatively few, comparatively large, comparatively predictable orders. A warehouse
serving e-commerce handles many small orders with heterogeneous urgency, several service tiers,
and pronounced seasonal peaks.

Their more recent retrospective across five decades of warehousing research
\citep{boysen2025} reinforces the same point: the operational questions have moved faster than
the modelling apparatus, and workforce-related decisions in particular remain comparatively
under-served relative to their practical weight.

Three implications of this shift matter for the present study. Demand is
\textbf{non-stationary}, so any model assuming steady state will mis-describe the peaks where
plans actually fail. Orders are \textbf{heterogeneous}, so an order is not a reliable unit of
work. Service commitments are \textbf{differentiated}, so a single aggregate performance measure
is not sufficient to determine whether a plan is acceptable.

\subsection{Order heterogeneity and stage-differentiated workload}

The order-picking literature has long recognised that orders differ in the work they require:
order size, item characteristics and storage location all affect picking time
\citep{dekoster2007,chan2011}. A related strand argues that even this understates the variation,
because human factors --- learning, fatigue, task variety --- introduce systematic differences
that deterministic models omit \citep{grosse2014}.

Design-oriented studies formalise this variation into system-selection guidance.
\citet{dallari2008} propose a classification of order-picking system designs and the conditions
favouring each, while \citet{marchet2011} develop an analytical model for estimating the
performance of pick-and-sort systems at the design stage. Complementary case-based work
\citep{marchet2014} examines how such systems are chosen in practice and finds that the decision
is driven as much by order-profile characteristics --- order size, urgency mix, volume
variability --- as by throughput requirements alone. These studies reinforce the point that
order composition, not just order count, determines what an operation must be sized to handle.

Less attention has been paid to a further point that is central here: the same order attribute
can affect different stages by different amounts. A fragile item may be only slightly slower to
retrieve but substantially slower to pack safely; a bulky item may be the reverse. Where studies
model multiple warehouse stages, workload is frequently treated as proportional across them, so
that an order demanding at one stage is assumed demanding at all. If that assumption fails, the
consequence is not a small error but a mis-allocation of capacity between stages --- and, as
\Cref{sec:res-demand} shows, in the operation studied here the ratio of workload across stages
is roughly \(1 : 0.72 : 0.34\) and shifts seasonally.

\section{Capacity and workforce planning}
\label{sec:lit-capacity}

\subsection{From workload to headcount}

The conventional route from demand to headcount is a ratio: expected workload divided by
capacity per worker, adjusted by a target utilisation and rounded up. Its appeal is that it is
transparent, fast and requires only aggregate data. Its weakness is that it treats the
production system as a reservoir rather than a queue.

Queueing theory explains why that matters. In any system with variability in arrivals or
service, waiting time grows non-linearly with utilisation and diverges as utilisation approaches
one; the practical consequence, developed at length by \citet{hopp2011}, is that a system sized
to its average capacity requirement will not perform to its average expectations. This is the
reason that a utilisation target below one appears in the analytical estimate used in this
thesis, and it establishes the direction of the correction that queueing imposes on ratio-based
sizing.

The relationship that connects these quantities is Little's law \citep{little1961}, which
states that the average number of items in a stable system equals the arrival rate multiplied by
the average time an item spends there. Its relevance to the present problem is more than
notational. Throughput over a fixed horizon is bounded by capacity, so a workforce that cannot
process the month's arrivals will accumulate work-in-process regardless of how that work is
ordered. Sequencing does not create worker-minutes, and cannot reduce the aggregate the law
relates.

The boundary of that argument should be marked, because it is easy to over-extend. Little's law
is a statement about long-run averages in a stable system; it does not by itself imply that the
\emph{number} of orders completed before a finite cut-off is the same under every sequencing
discipline. Where service times are heterogeneous, as they are here, ordering can shift a few
completions across the horizon boundary. What the law supports is the direction and the
magnitude ceiling of the effect, not an identity. \Cref{sec:res-backlog} reports the
corresponding empirical result: at the configurations tested, backlog differed across policies
by fractions of a percentage point.

A related tradition in production control makes the same point from the opposite direction.
Workload control approaches manage congestion by regulating the release of work into the shop
rather than by adding capacity, on the argument that controlling what enters the system is often
cheaper than expanding it \citep{thurer2012}. The parallel with this thesis is close: both treat
the ordering and timing of work as a control lever operating on a fixed capacity base. The
difference is that workload control decides \emph{when} work enters, whereas the sequencing
policies here decide \emph{which} waiting work is served next, with capacity itself as the
second decision variable.

Stochastic modelling of warehouse operations is itself an established field, and
\citet{gong2011} review the models available and the operational questions they can address.
Their survey is instructive for the present purpose partly because of what it reveals about
scope: analytically tractable models generally treat one stage, one customer class, or
stationary demand, and the tractability is obtained precisely by not representing the
combination of features that makes tactical capacity planning difficult.

What queueing theory does \emph{not} supply, for a system of the kind studied here, is a closed
form. The classical results depend on assumptions --- stationary arrivals, exponential or
otherwise tractable service distributions, a single customer class or a simple priority
structure --- that are individually reasonable and jointly incompatible with a warehouse serving
seasonal, bursty, heterogeneous, multi-class demand across several stages in series. The
literature accordingly turns to simulation when these features must be retained simultaneously
\citep{gu2010,gong2011}.

\subsection{Workforce planning specifically}

Research addressing warehouse workforce decisions directly is comparatively sparse relative to
the picking-efficiency literature. \citet{ganbold2020} combine discrete-event simulation with a
neighbourhood search to assign warehouse workers, establishing simulation-based optimisation as
a workable approach to the staffing question and providing the closest methodological analogue
to the capacity component of this thesis. Related work in warehouse capacity design has combined
simulation with structured multi-criteria evaluation to support facility-level decisions
\citep{sencer2021}.

Both belong to the broader tradition of simulation optimisation, in which a simulation model
serves as the objective function for a search over decision variables that cannot be optimised
analytically. \citet{amaran2016} review the algorithms available and the conditions under which
each is appropriate, and their framing clarifies a choice made in this thesis. Where the
decision space is large and the simulation expensive, sophisticated search is warranted; where
the decision space is a small set of operationally plausible configurations, exhaustive
evaluation of a bounded candidate set is both simpler and more transparent. The present study
takes the latter route, and consequently makes no optimality claim.

Two features distinguish the problem addressed here from that body of work. The first is the
planning unit: this thesis works in monthly full-time equivalents, a tactical capacity quantity,
rather than in task assignments or facility dimensions. The second, and more important, is that
the workforce decision is not taken in isolation from the operating discipline. In most
workforce studies the sequencing rule is a fixed background condition; here it is a second
decision variable, and the interaction between the two is the object of study.

\section{Sequencing, priority and differentiated service}
\label{sec:lit-sequencing}

\subsection{Dispatching rules}

Deciding which waiting job to process next is among the oldest problems in production research,
and the catalogue of rules proposed for it is correspondingly large. The survey by
\citet{panwalkar1977} remains the canonical organisation of that catalogue, classifying rules by
the information they use --- arrival time, processing time, due date, or combinations --- and
summarising their behaviour against standard objectives.

Two findings from that tradition are directly relevant. First, no rule dominates across all
objectives and conditions; the appropriate rule depends on what is being optimised and under
what load. Second, when jobs carry due dates, rules that use due-date information systematically
outperform rules that do not, and the advantage widens as the system becomes more congested.
First-in-first-out, which uses only arrival order, is consequently the natural baseline against
which due-date-aware rules are measured --- and the natural expectation is that it will perform
poorly when deadlines differ, which is precisely what \Cref{sec:res-policy} observes.

\subsection{Priority queues and multi-class service}

Where customers fall into classes with different requirements, the standard response is a
priority discipline. Classical priority queueing, developed from the early 1960s
\citep{aviitzhak1961}, establishes the essential trade-off: raising a class's priority reduces
its waiting time at the direct expense of lower-priority classes, and under a non-pre-emptive
discipline the total work in the system is unchanged --- priority redistributes waiting rather
than removing it. The distinction between pre-emptive and non-pre-emptive service matters for
what can be promised: under a non-pre-emptive rule, an arriving urgent order must still wait for
whatever is currently in service, which places a floor on its achievable response time that no
priority setting can remove. The engines studied here are non-pre-emptive throughout, so that
floor applies to all three policies equally.

That conservation property is the theoretical foundation of the central empirical result of this
thesis. If strict priority merely moves waiting from the urgent class to the normal class, then
where the normal class also carries a binding commitment, there must exist conditions under
which protecting the urgent class pushes the normal class below its own requirement. Strict
priority then fails, not because it performs its function badly, but because it performs it
without limit. \Cref{sec:disc-policy} reports precisely this failure in the peak month.

The conservation argument also suggests something that is easy to test and that
\Cref{sec:res-backlog} bears out: if a non-pre-emptive discipline mainly redistributes waiting,
then the total amount of work completed within a fixed horizon should be close to independent of
the discipline chosen. The qualifier is deliberate. Conservation constrains the \emph{work}
completed, expressed in worker-minutes, exactly; it constrains the \emph{count} of orders
completed before a finite cut-off only approximately, because service times are heterogeneous
and a discipline that favours shorter jobs can push a few extra completions inside the horizon.
The prediction to test is therefore near-invariance rather than identity, and near-invariance is
what the results show. This provides a useful internal consistency check on the model as well as
a boundary on what any sequencing policy can be expected to achieve --- a policy cannot
manufacture capacity, and where a workforce is simply too small, no ordering of the queue will
make it sufficient.

A related implication concerns the direction in which the trade-off can be pushed. Because the
waiting removed from one class must appear somewhere, a policy that is free to moderate its own
priority --- serving the lower-priority class when the higher-priority class has sufficient
slack --- occupies a strictly larger space of outcomes than one that cannot. Whether that
additional freedom is worth anything depends on whether the constraint on the lower-priority
class is ever close to binding, which is an empirical question about the operating regime rather
than a theoretical one.

A subtlety that receives less attention in applied work concerns ordering \emph{within} a
priority class. Theoretical treatments generally assume first-in-first-out within class, which
keeps the waiting-time distribution of each class comparatively concentrated. Implementations
using heap-based priority structures do not necessarily provide that guarantee, and where the
performance criterion is a threshold rather than a mean, the tail of the distribution is what
determines the outcome. \Cref{sec:urgent-first} documents exactly this situation in the baseline
studied here.

\subsection{Differentiated service levels}

The practice of committing to different service levels for different customer segments is
widespread, and it changes the structure of the planning problem in a way that is easy to
overlook. When performance is measured in aggregate --- mean tardiness, overall on-time
percentage, makespan --- the objective is a single scalar and improvements anywhere are
fungible. When commitments are class-specific and contractual, the requirement becomes a
\emph{conjunction} of constraints, and improvement in one class cannot compensate for failure in
another.

This distinction has a direct consequence for evaluation. A configuration may score well on any
aggregate measure while failing a minority class completely, because the minority's failures are
diluted by the majority's success. \Cref{sec:res-policy} provides a concrete instance: a
configuration attaining \SI{81}{\percent} overall service delivered \SI{24.5}{\percent} to its
urgent customers. Much of the dispatching-rule literature evaluates against aggregate
objectives, under which a strict priority rule appears uniformly favourable; under a
conjunctive class-specific constraint it does not.

\section{Discrete-event simulation}
\label{sec:lit-des}

\subsection{Why the method fits the problem}

Discrete-event simulation represents a system as a sequence of instantaneous events that change
its state, advancing the clock from one event to the next \citep{law2015}. It is the standard
approach where the quantities of interest --- queue lengths, waiting-time distributions,
resource utilisation, congestion propagation --- emerge from interactions between individual
entities and finite resources rather than from a tractable aggregate relationship.

The features that make analytical treatment difficult are exactly those that discrete-event
simulation accommodates without special provision: arrivals may follow any pattern, including
the bursty and seasonal patterns observed in fulfilment; service requirements may vary by entity
and by stage; resources may be finite and multiple; and the queue discipline may be arbitrary,
including a discipline determined at run time by a learned policy. The cost is that results are
statistical estimates from a model rather than closed-form properties of a system, which places
weight on experimental design and on validation.

\subsection{Model credibility}

Because a simulation produces numbers regardless of whether it represents anything, credibility
must be established explicitly. \citet{sargent2008} provides the framework used in this thesis,
distinguishing \emph{verification} --- whether the computerised model implements the conceptual
model correctly --- from \emph{validation} --- whether the conceptual model is an adequate
representation for the intended purpose. The distinction matters practically: a model can be
correctly implemented and invalid, or approximately implemented and adequate.

\citet{robinson2008} adds the complementary point that a simulation model's fitness is
determined largely at the conceptual modelling stage, by decisions about scope and level of
detail taken before any code is written. Those decisions --- here, to represent capacity and
queueing while excluding physical warehouse geography --- define what questions the model can
answer and are the primary source of its limitations.

\citet{kleijnen1995} sharpens what validation can realistically mean when, as here, no
real-system data are available to compare against. In that situation the achievable checks are
internal: that the model reproduces behaviours theory predicts, that it responds sensibly to
changes in its inputs, and that its assumptions are stated explicitly enough for a reader to
judge them. Validation in that sense establishes a model as a coherent instrument for comparing
alternatives under stated assumptions, not as a predictor of a particular real system --- which
is precisely the claim made in \Cref{sec:verification} and no more.

This distinction also determines what kind of research design the study belongs to.
\citet{bertrand2002} classify quantitative modelling research in operations management along two
axes: whether the model is axiomatic (studying properties of an idealised model) or empirical
(seeking to match observed reality), and whether the work is descriptive or normative. On that
classification the present study is \emph{axiomatic normative}: it builds a model, defines
decision variables and a performance criterion, and derives recommendations within the model's
assumptions, without claiming that the model has been fitted to a real operation. Recognising
this explicitly guards against a common overreach --- reporting model outcomes as though they
were empirical findings about warehouses in general --- and it is the reason this thesis is
framed as applied simulation-based experimental research rather than as an empirical study.

\section{Simulation experimentation}
\label{sec:lit-experimentation}

A simulation model is an instrument; obtaining reliable conclusions from it is a matter of
experimental design. Two considerations are central to this thesis.

The first is \textbf{comparison under common conditions}. When alternative system
configurations are compared, differences in the output may reflect genuine differences between
configurations or merely different draws from the same random processes. Common random numbers
address this by ensuring that alternatives face identical stochastic realisations, which
induces positive correlation between their outputs and reduces the variance of the estimated
\emph{difference} \citep{law2015}. It is the standard variance reduction technique for
comparing systems, and it is what allows the policy comparison in this thesis to be made
meaningfully without large replication counts.

Applying it correctly requires care, and the naive implementation fails. If service times are
drawn as entities enter service, then policies that reorder the queue consume the random stream
in a different sequence, and a shared seed no longer produces shared service times. The
synchronisation must be at the level of the entity and the operation, not the draw order ---
which is why the framework described in \Cref{sec:crn} pre-computes service times per (order,
stage) before any policy runs.

The second consideration is \textbf{what a comparison under common random numbers does not
establish}. It makes the difference between policies precise at a given demand realisation; it
says nothing about how that difference would vary across realisations. Establishing the latter
requires replication, and the asymmetry between the two workflows in this thesis --- one
realisation per month retrospectively, three replications prospectively --- reflects that
distinction rather than an oversight.

\section{Reinforcement learning for operational control}
\label{sec:lit-rl}

\subsection{Foundations}

Reinforcement learning addresses sequential decision problems formulated as Markov decision
processes: an agent observes a state, selects an action, receives a reward and transitions to a
new state, and learns a policy maximising cumulative reward \citep{suttonbarto2018}. Q-learning
estimates the value of taking an action in a state and following the optimal policy thereafter,
and does so without a model of the environment's dynamics --- an attraction where those dynamics
are complex, as in a congested multi-stage queueing system.

The extension that made the approach practical for large state spaces was the deep Q-network of
\citet{mnih2015}, which approximates the action-value function with a neural network and
stabilises training through experience replay and a periodically updated target network. That
combination --- function approximation, replay and target network --- is the algorithm employed
by the learned policy evaluated in this thesis.

\subsection{Application to scheduling and operations}

Applying reinforcement learning to production scheduling has become an active area.
\citet{luo2020} is the closest analogue to the present work: a deep Q-network selects among
dispatching decisions in a flexible job shop subject to dynamic job arrivals, a problem sharing
the essential structure addressed here --- repeated sequencing decisions under uncertainty about
what will arrive next. The approach has since extended into operations management more broadly,
including multi-echelon inventory control \citep{liu2024}, indicating that learned controllers
are being taken seriously in mainstream operational research rather than only in machine
learning venues. Within warehousing specifically, \citet{mahmoudinazlou2025} apply deep
reinforcement learning to dynamic order picking, confirming that the technique is now being
brought to bear on warehouse operational control rather than only on manufacturing scheduling.

A cautionary strand accompanies this growth. \citet{waubert2022} examine the reliability of
reinforcement learning based scheduling systems and note that reported performance can depend
heavily on the conditions under which an agent was trained, so that behaviour outside those
conditions is not guaranteed. That concern is directly relevant here for two reasons. It
motivates the decision to hold a subset of workforce configurations out of training entirely, so
that the agent's behaviour on unseen capacity can be observed. And it argues for reporting the
regime in which a learned policy helps rather than an aggregate performance figure, since an
average conceals precisely the conditional behaviour that reliability concerns are about.

Two observations about this literature shape how the present study frames its own contribution.
First, the learned policy is normally the object of study, evaluated against benchmark instances
and reported as a headline performance figure. Second, the objectives are typically aggregate
--- mean tardiness, makespan, total cost --- rather than conjunctive class-specific constraints.
Neither is a criticism, but both mean that the question ``under what conditions does the learned
policy help, and by how much?'' is asked less often than ``does it win?''. This thesis is
positioned to ask the former, because the learned policy here is one component among several
inside a capacity-planning framework rather than the subject of the investigation.

\section{Decision support and explainability}
\label{sec:lit-dss}

A model that produces a correct recommendation which nobody adopts has not supported a decision.
The point is long-standing in the management-science literature. \citet{little1970}, introducing
the notion of a \emph{decision calculus}, argued that a model intended for managerial use must
be simple enough to be understood, robust, easy to control and complete on important issues, and
observed that models failing these criteria are not used regardless of their technical quality.
That argument treats interpretability as a functional requirement of the instrument rather than
as presentation, which is the position taken here, and it matters particularly where a
recommendation carries budgetary consequences a manager must defend.

Two requirements follow for a capacity-planning instrument. It must expose \emph{why} it
recommends what it does --- which stage constrains the operation, which service class is at
risk, what drives the cost --- because a recommendation that cannot be interrogated cannot be
checked or corrected. And it must express its recommendations in the units in which decisions
are actually taken: headcount, cost, and how often a commitment is met across the scenarios
evaluated, rather than utilisation ratios or abstract objective values.

This requirement interacts directly with the choice of sequencing policy. A learned policy is
opaque in a way that a two-line priority rule is not, and \citet{rudin2019} argues that in
high-stakes decisions this opacity is not neutral: a post-hoc explanation of a black-box model is
not the model's actual reasoning, so an interpretable model should be preferred unless the
opaque alternative demonstrably performs better. Applied here, that sets the terms on which a
learned sequencing policy can be recommended at all --- the size and breadth of its advantage
over a transparent heuristic must be established rather than assumed, which is precisely what
\Cref{sec:res-policy} sets out to measure. \Cref{sec:disc-literature} returns to the theme.

\section{Positioning of this thesis}
\label{sec:lit-gap}

The literature reviewed above establishes each of the components this thesis uses. Warehouse
operations research characterises the setting and the dominance of picking cost. Queueing theory
explains why ratio-based capacity sizing is unreliable under variability. Scheduling research
catalogues dispatching rules and establishes that due-date-aware rules outperform
first-in-first-out under congestion. Priority queueing establishes that protecting one class
redistributes waiting to another. Discrete-event simulation provides the instrument for studying
systems where these features coexist, and common random numbers provide the means of comparing
alternatives fairly. Reinforcement learning provides a means of deriving a sequencing policy
rather than specifying one.

What this thesis contributes is not a new component but their \emph{integration} within a single
decision-support framework, and the empirical characterisation that integration makes possible.
No claim is made that the individual elements are novel, nor that no prior study has combined
some of them; the field is large and such a claim could not be substantiated.

The specific position occupied is the following. Existing work generally treats capacity and
sequencing as separate problems, evaluates sequencing rules against aggregate objectives, and
reports whether a method wins rather than where. This thesis treats capacity and sequencing as
one joint decision, evaluates it against a conjunction of class-specific service constraints
together with an explicit economic objective, compares a learned policy against transparent
heuristics on identical stochastic realisations, and pairs the recommendation with a bottleneck
diagnostic and an explicit economic test of marginal capacity.

That position generates the four research questions stated in \Cref{sec:research-questions}, and
it is what allows the study to produce findings that a component-wise treatment would not
surface: that ratio-based sizing over-provides capacity because it cannot represent the
substitution between capacity and operating discipline; that the value of a sequencing policy is
conditional on the capacity regime; and that a correctly identified bottleneck is frequently not
an economically justified place to add capacity.
